\documentclass[12pt]{article}
\usepackage[margin=0.75in]{geometry}
\usepackage{fontspec}
\setmainfont{Latin Modern Roman}
\usepackage{microtype}
\usepackage{multicol}
\usepackage{fancyhdr}
\usepackage{titlesec}
\usepackage{enumitem}
\usepackage{xcolor}
\usepackage[hidelinks]{hyperref}
\setlength{\columnsep}{0.28in}
\setlength{\parindent}{1.1em}
\setlength{\parskip}{0.32em}
\setlength{\headheight}{15pt}
\titleformat{\section}{\large\bfseries\scshape}{}{0pt}{}
\titleformat{\subsection}{\normalsize\bfseries}{}{0pt}{}
\pagestyle{fancy}
\fancyhf{}
\fancyfoot[C]{\thepage}
\renewcommand{\headrulewidth}{0.4pt}
\newcommand{\focusbox}[1]{\par\vspace{0.4em}\noindent\begingroup\setlength{\fboxsep}{6pt}\colorbox{gray!12}{\begin{minipage}{\dimexpr\linewidth-2\fboxsep\relax}#1\end{minipage}}\endgroup\par\vspace{0.5em}}

\fancyhead[L]{Whately: Logic and Rhetoric}
\fancyhead[R]{Student Reading}
\begin{document}
\begin{center}
{\LARGE\bfseries Persuasion Is Not Proof}\par\vspace{0.25em}
{\large\scshape Week 16: Logic and Rhetoric}\par\vspace{0.5em}
{Adapted and modernized from Richard Whately, \textit{Elements of Logic}}\par\vspace{0.6em}
\rule{0.82\textwidth}{0.4pt}
\end{center}

\section*{Central Question}
How do style, emotion, character, and persuasion relate to proof?

\section*{Vocabulary Preview}
\begin{multicols}{2}
\begin{itemize}[leftmargin=*, itemsep=0.18em]
\item \textbf{Rhetoric}: the art of speaking or writing to move an audience.
\item \textbf{Proof}: reasons and evidence that actually support a claim.
\item \textbf{Persuasion}: what makes an audience believe, feel, or act.
\item \textbf{Logos}: appeal through reasons, structure, and evidence.
\item \textbf{Ethos}: appeal through the speaker's character, credibility, or standing.
\item \textbf{Pathos}: appeal through emotion, vivid scene, or feeling.
\item \textbf{Audience}: the people the speaker is trying to move.
\item \textbf{Rhetorical force}: how powerfully a speech lands, regardless of proof.
\item \textbf{Claim under proof}: the exact conclusion that still needs support.
\item \textbf{Rhetoric-vs-proof analysis}: separate what is felt, what is claimed, and what is proved.
\end{itemize}
\end{multicols}

\section*{Introduction}
A speech can move a room and still fail as proof. Another speech can be carefully reasoned and still fail to persuade a tired or hostile audience. Whately's larger point about the province of logic matters here: logic tests whether reasons support a conclusion. Rhetoric asks how language, character, and emotion reach hearers. Both matter. They are not the same job.

Students meet this distinction constantly. A coach's pep talk, a campaign ad, a funeral oration, a social-media post, and a classroom debate may all be powerful without being sound. The disciplined habit is not to despise rhetoric, and not to confuse it with proof.

\focusbox{\textbf{Source note:} Adapted and modernized from Whately's treatment of the province of logic, with classroom terms from classical rhetoric (\textit{logos}, \textit{ethos}, \textit{pathos}). The goal is not to flatten persuasion into ``manipulation,'' but to separate rhetorical force from logical support. Classroom examples are original.}

\begin{multicols}{2}
\section*{Reading}

\subsection*{1. Two Different Questions}
When you hear a strong speech, ask two questions:

\begin{itemize}[leftmargin=*]
\item What does this make the audience feel or want to do?
\item What claim, if any, does it actually prove?
\end{itemize}

Those questions can have different answers. A vivid story may make listeners angry. That does not yet prove the speaker's policy is best. A calm list of evidence may prove a limited claim. That does not yet move a crowd.

\subsection*{2. Logos, Ethos, and Pathos}
Classical rhetoric names three common appeals:

\begin{itemize}[leftmargin=*]
\item \textbf{Logos}: reasons, comparisons, examples, and structure.
\item \textbf{Ethos}: the speaker's honesty, expertise, loyalty, or public standing.
\item \textbf{Pathos}: pity, fear, pride, shame, grief, or hope.
\end{itemize}

None of these is automatically dishonest. A true claim may need clear reasons, a trustworthy speaker, and language the audience can feel. The error is treating any one appeal as a substitute for the burden of proof.

\subsection*{3. Rhetorical Force Is Not Validity}
A speech can be \textit{effective} without being \textit{sound}. Effectiveness is about audience movement: tears, applause, votes, action. Soundness still requires true or well-supported premises and a conclusion that follows.

This is why Week 6 form-versus-fact work and Week 10 missing-the-point work still matter. A speaker may prove a nearby claim, stir emotion about a different claim, or use character praise to hide a missing middle step.

\subsection*{4. What Rhetoric Can Honestly Do}
Rhetoric can help truth become visible. Clear arrangement helps hearers follow proof. Credible ethos can rightly increase attention. Pathos can make a real harm feel concrete rather than abstract. Rhetoric can also expose a weak claim by making its costs vivid.

So the disciplined student does not say, ``Rhetoric is fake.'' The student says, ``Show me what is claimed, what is felt, and what is proved.''

\subsection*{5. What Rhetoric Cannot Supply by Itself}
Rhetoric cannot invent missing facts. It cannot turn a circular premise into independent support. It cannot make an irrelevant story answer the point at issue. It cannot convert a partial example into a universal claim merely because the story is moving.

If the audience leaves convinced, ask: convinced of what, and on what grounds?

\subsection*{6. How to Run a Rhetoric-vs-Proof Analysis}
Use this map:

\begin{enumerate}[leftmargin=*,itemsep=0.12em]
\item State the speaker's main claim.
\item Name the audience and the desired response.
\item Identify major appeals: logos, ethos, pathos.
\item Separate emotional or character force from stated reasons.
\item Ask what has been proved, what has only been suggested, and what still needs evidence.
\item Judge whether the rhetorical force matches, exceeds, or replaces the proof.
\end{enumerate}

\subsection*{7. Literature and Public Speech}
In \textit{Julius Caesar}, Antony's funeral speech is the classic case: grief, irony, personal tokens, and public loyalty reshape the crowd. In \textit{Macbeth}, Lady Macbeth uses shame, manhood, and vivid courage-language to press Macbeth toward action. Both passages reward a rhetoric-vs-proof analysis rather than a simple label of ``good speech'' or ``bad manipulation.''

\section*{Reading Focus}
\begin{enumerate}[leftmargin=*]
\item What is the difference between persuasion and proof?
\item Define logos, ethos, and pathos in one sentence each.
\item Why can a speech be effective without being sound?
\item What honest work can rhetoric do for truth?
\item What can rhetoric not supply by itself?
\item Copy the six steps of the rhetoric-vs-proof analysis.
\item Name one speech from \textit{Macbeth} or \textit{Julius Caesar} that is strong rhetorically and still needs logical testing.
\end{enumerate}

\section*{Handout Summary}
Write 5--7 sentences explaining how persuasion differs from proof. Your summary must include the words \textit{logos}, \textit{ethos}, \textit{pathos}, \textit{rhetorical force}, and \textit{claim under proof}.
\end{multicols}
\end{document}
